\section{Serre subcategories}
\label{section-serre-subcategories}

\noindent
In \cite[Chapter I, Section 1]{Serre_homotopie_classes}
a notion of a ``class'' of abelian groups is defined.
This notion has been extended to abelian categories by many authors
(in slightly different ways). We will use the following variant
which is virtually identical to Serre's original definition.

\begin{definition}
\label{definition-serre-subcategory}
\begin{reference}
\cite[Condition (I) on page 259]{Serre_homotopie_classes}
\end{reference}
Let $\mathcal{A}$ be an abelian category.
\begin{enumerate}
\item A {\it Serre subcategory} of $\mathcal{A}$ is a
nonempty full subcategory $\mathcal{C}$ of $\mathcal{A}$
such that given an exact sequence\footnote{By
Definition \ref{definition-exact} this means $\Im(A \to B) = \Ker(B \to C)$.}
$$
A \to B \to C
$$
with $A, C \in \Ob(\mathcal{C})$, then also
$B \in \Ob(\mathcal{C})$.
\item A {\it weak Serre subcategory} of $\mathcal{A}$ is a nonempty
full subcategory $\mathcal{C}$ of $\mathcal{A}$ such that given an
exact sequence
$$
A_0 \to A_1 \to A_2 \to A_3 \to A_4
$$
with $A_0, A_1, A_3, A_4$ in $\mathcal{C}$, then also $A_2$ in $\mathcal{C}$.
\end{enumerate}
\end{definition}

\noindent
In some references the second notion is called a ``thick'' subcategory
and in other references the first notion is called a ``thick'' subcategory.
However, it seems that the notion of a Serre subcategory is universally
accepted to be the one defined above. Note that in both cases the category
$\mathcal{C}$ is abelian and that the inclusion functor
$\mathcal{C} \to \mathcal{A}$ is a fully faithful exact functor.
Let's characterize these types of subcategories in more detail.

\begin{lemma}
\label{lemma-characterize-serre-subcategory}
Let $\mathcal{A}$ be an abelian category.
Let $\mathcal{C}$ be a subcategory of $\mathcal{A}$.
Then $\mathcal{C}$ is a Serre subcategory if and only if
the following conditions are satisfied:
\begin{enumerate}
\item $0 \in \Ob(\mathcal{C})$,
\item $\mathcal{C}$ is a strictly full subcategory of $\mathcal{A}$,
\item any subobject or quotient of an object of $\mathcal{C}$ is an object
of $\mathcal{C}$,
\item if $A \in \Ob(\mathcal{A})$ is an extension of objects of $\mathcal{C}$
then also $A \in \Ob(\mathcal{C})$.
\end{enumerate}
Moreover, a Serre subcategory is an abelian category and
the inclusion functor is exact.
\end{lemma}

\begin{proof}
Omitted.
\end{proof}

\begin{lemma}
\label{lemma-characterize-weak-serre-subcategory}
Let $\mathcal{A}$ be an abelian category.
Let $\mathcal{C}$ be a subcategory of $\mathcal{A}$.
Then $\mathcal{C}$ is a weak Serre subcategory if and only if
the following conditions are satisfied:
\begin{enumerate}
\item $0 \in \Ob(\mathcal{C})$,
\item $\mathcal{C}$ is a strictly full subcategory of $\mathcal{A}$,
\item kernels and cokernels in $\mathcal{A}$ of morphisms
between objects of $\mathcal{C}$ are in $\mathcal{C}$,
\item if $A \in \Ob(\mathcal{A})$ is an extension of objects of $\mathcal{C}$
then also $A \in \Ob(\mathcal{C})$.
\end{enumerate}
Moreover, a weak Serre subcategory is an abelian category and
the inclusion functor is exact.
\end{lemma}

\begin{proof}
Omitted.
\end{proof}

\begin{lemma}
\label{lemma-kernel-exact-functor}
Let $\mathcal{A}$, $\mathcal{B}$ be abelian categories.
Let $F : \mathcal{A} \to \mathcal{B}$ be an exact functor.
Then the full subcategory of objects $C$ of $\mathcal{A}$
such that $F(C) = 0$ forms a Serre subcategory of $\mathcal{A}$.
\end{lemma}

\begin{proof}
Omitted.
\end{proof}

\begin{definition}
\label{definition-kernel-category}
Let $\mathcal{A}$, $\mathcal{B}$ be abelian categories.
Let $F : \mathcal{A} \to \mathcal{B}$ be an exact functor.
Then the full subcategory of objects $C$ of $\mathcal{A}$
such that $F(C) = 0$ is called the {\it kernel of the functor $F$},
and is sometimes denoted $\Ker(F)$.
\end{definition}

\noindent
Any Serre subcategory of an abelian category is the kernel of
an exact functor. In
Examples, Section \ref{examples-section-serre-quotient-modulo-torsion-modules}
we discuss this for Serre's original example of torsion groups.

\begin{lemma}
\label{lemma-serre-subcategory-is-kernel}
Let $\mathcal{A}$ be an abelian category.
Let $\mathcal{C} \subset \mathcal{A}$ be a Serre subcategory.
There exists an abelian category $\mathcal{A}/\mathcal{C}$
and an exact functor
$$
F : \mathcal{A} \longrightarrow \mathcal{A}/\mathcal{C}
$$
which is essentially surjective and whose kernel is $\mathcal{C}$.
The category $\mathcal{A}/\mathcal{C}$ and the functor $F$ are
characterized by the following universal property: For any exact
functor $G : \mathcal{A} \to \mathcal{B}$ such that
$\mathcal{C} \subset \Ker(G)$ there exists a factorization
$G = H \circ F$ for a unique exact functor
$H : \mathcal{A}/\mathcal{C} \to \mathcal{B}$.
\end{lemma}

\begin{proof}
Consider the set of arrows of $\mathcal{A}$ defined by
the following formula
$$
S = \{f \in \text{Arrows}(\mathcal{A}) \mid
\Ker(f), \Coker(f) \in \Ob(\mathcal{C}) \}.
$$
We claim that $S$ is a multiplicative system. To prove this we have
to check MS1, MS2, MS3, see
Categories, Definition \ref{categories-definition-multiplicative-system}.

\medskip\noindent
It is clear that identities are elements of $S$. Suppose that
$f : A \to B$ and $g : B \to C$ are elements of $S$.
There are exact sequences
$$
\begin{matrix}
0 \to \Ker(f) \to \Ker(gf) \to \Ker(g) \\
\Coker(f) \to \Coker(gf) \to \Coker(g) \to 0
\end{matrix}
$$
Hence it follows that $gf \in S$. This proves MS1. (In fact, a similar
argument will show that $S$ is a saturated multiplicative system, see
Categories, Definition
\ref{categories-definition-saturated-multiplicative-system}.)

\medskip\noindent
Consider a solid diagram
$$
\xymatrix{
A \ar[d]_t \ar[r]_g & B \ar@{..>}[d]^s \\
C \ar@{..>}[r]^f & C \amalg_A B
}
$$
with $t \in S$. Set
$W = C \amalg_A B =  \Coker((t, -g) : A \to C \oplus B)$.
Then $\Ker(t) \to \Ker(s)$ is surjective and
$\Coker(t) \to \Coker(s)$ is an isomorphism. Hence
$s$ is an element of $S$. This proves LMS2 and the proof of RMS2 is dual.

\medskip\noindent
Finally, consider morphisms $f, g : B \to C$ and a morphism $s : A \to B$
in $S$ such that $f \circ s = g \circ s$. This means that
$(f - g) \circ s = 0$. In turn this means that
$I = \Im(f - g) \subset C$ is a quotient of $\Coker(s)$
hence an object of $\mathcal{C}$. Thus $t : C \to C' = C/I$ is an
element of $S$ such that $t \circ (f - g) = 0$, i.e., such that
$t \circ f = t \circ g$. This proves LMS3 and the proof of
RMS3 is dual.

\medskip\noindent
Having proved that $S$ is a multiplicative system we set
$\mathcal{A}/\mathcal{C} = S^{-1}\mathcal{A}$, and we set
$F$ equal to the localization functor $Q$. By
Lemma \ref{lemma-localization-abelian}
the category $\mathcal{A}/\mathcal{C}$ is abelian and $F$ is exact.
If $X$ is in the kernel of $F = Q$, then by
Lemma \ref{lemma-kernel-localization}
we see that $0 : X \to Z$ is an element of $S$ and hence
$X$ is an object of $\mathcal{C}$, i.e., the kernel of
$F$ is $\mathcal{C}$.
Finally, if $G$ is as in the statement of the lemma, then $G$ turns
every element of $S$ into an isomorphism. Hence we obtain the
functor $H : \mathcal{A}/\mathcal{C} \to \mathcal{B}$ from
the universal property of localization, see
Categories, Lemma \ref{categories-lemma-properties-left-localization}.
We still have to show the functor $H$ is exact.
To do this it suffices to show that $H$ commutes
with taking kernels and cokernels, see Lemma \ref{lemma-exact-functor}.
Let $A \to B$ be a morphism in $\mathcal{A}/\mathcal{C}$.
We may represent $A \to B$ as $fs^{-1}$ where $s : A' \to A$
is in $S$ and $f : A' \to B$ an arbitrary morphism of $\mathcal{A}$.
Since $F = Q$ maps $s$ to an isomorphism in the quotient category
$\mathcal{A}/\mathcal{C}$, it suffices to show that $H$ commutes with taking
kernels and cokernels of morphisms $f : A \to B$ of $\mathcal{A}$.
But here we have $H(f) = G(f)$ and the result follows
from the fact that $G$ is exact.
\end{proof}

\begin{lemma}
\label{lemma-quotient-by-kernel-exact-functor}
Let $\mathcal{A}$, $\mathcal{B}$ be abelian categories.
Let $F : \mathcal{A} \to \mathcal{B}$ be an exact functor.
Let $\mathcal{C} \subset \mathcal{A}$ be a Serre subcategory
contained in the kernel of $F$.
Then $\mathcal{C} = \Ker(F)$ if and only if the induced functor
$\overline{F} : \mathcal{A}/\mathcal{C} \to \mathcal{B}$
(Lemma \ref{lemma-serre-subcategory-is-kernel}) is faithful.
\end{lemma}

\begin{proof}
We will use the results of Lemma \ref{lemma-serre-subcategory-is-kernel}
without further mention.
The ``only if'' direction is true because the kernel of $\overline{F}$ is zero
by construction. Namely, if $f : X \to Y$ is a morphism in
$\mathcal{A}/\mathcal{C}$ such that $\overline{F}(f) = 0$, then
$\overline{F}(\Im(f)) = \Im(\overline{F}(f)) = 0$, hence $\Im(f) = 0$ by the
assumption on the kernel of $F$. Thus $f = 0$.

\medskip\noindent
For the ``if'' direction, let $X$ be an object of $\mathcal{A}$ such that $F(X)
= 0$. Then $\overline{F}(\text{id}_X) = \text{id}_{\overline{F}(X)} = 0$, thus
$\text{id}_X = 0$ in $\mathcal{A}/\mathcal{C}$ by faithfulness of
$\overline{F}$. Hence $X = 0$ in $\mathcal{A}/\mathcal{C}$, that is $X \in
\Ob(\mathcal{C})$.
\end{proof}






